\clearpage
\section{How to Run}

The mutual-exclusion behaviour can be observed manually by starting a local
RabbitMQ broker and then running two clients that target the same critical
section.

All commands below assume the current working directory is already
\texttt{assignment-4/distributed-critical-sections}.

\subsection{Broker startup}

Start RabbitMQ with:

\begin{lstlisting}[style=codestyle,language=bash,caption={Starting RabbitMQ from PowerShell}]
docker compose -f docker-compose.rabbitmq.yml up -d
\end{lstlisting}

\subsection{ProcessApp demo}

Open two PowerShell terminals and run the following commands:

\begin{lstlisting}[style=codestyle,language=bash,caption={Two overlapping ProcessApp clients from PowerShell}]
mvn exec:java "-Dexec.mainClass=pcd.dcs.demo.ProcessApp" "-Dexec.args=Process-A localhost 5672"
mvn exec:java "-Dexec.mainClass=pcd.dcs.demo.ProcessApp" "-Dexec.args=Process-B localhost 5672"
\end{lstlisting}

Both clients target the same critical-section name, \texttt{demo-cs}. Each
process performs multiple acquire/release cycles and keeps the token long
enough for the overlap to be visible from the console logs. One process enters
the critical section, while the other remains blocked in \texttt{enter()}
until the token is requeued. The console output and the shared file
\texttt{dcs\_shared.log} show alternating access without overlap.
